\section{The Angular Velocity Is Selected by the Dynamics}
\label{sec:simulation}

The rotation mechanism of Section~\ref{sec:rotation} yields a precise, falsifiable value for the
steady-state angular velocity. This section derives it as a selection theorem: \emph{if} the closed
loop settles onto a nondegenerate rigid rotation, its angular speed must be $\sqrt{k_2}$, whatever the
radii and whatever the shape. The theorem is confirmed numerically to $0.03\%$ over a factor of $20$ in
$k_2$; see Section~\ref{S-sec:sim_omega} of the simulation companion.

\subsection{A rigorous selection theorem}
In the target's frame let $\txi_i=x_i-x_0$, $\tvell_i=v_i-v_0$, and let
$J=\begin{bmatrix}0&-1\\1&0\end{bmatrix}$ denote the $90^\circ$ rotation matrix. The prediction
$|\omega|=\sqrt{k_2}$ follows from an exact necessary-condition argument. It does not assume a
regular polygon or local stability; it assumes only that a nondegenerate rigid rotation of the
closed-loop system exists.

\begin{proposition}[Single-integrator rotation-rate selection]\label{prop:omega}
Assume that, on a time interval of positive length, the closed loop
\eqref{eq:veh},\eqref{eq:tgt},\eqref{eq:ctrl7} admits a nondegenerate rigid rotation about the target,
\begin{equation}
  \txi_i(t)=R(\omega t)\,\txi_i^0,
  \qquad R(\theta)=\cos\theta\,I+\sin\theta\,J,
  \qquad \sum_{i\in\NN}\norm{\txi_i^0}^2>0 ,
  \label{eq:rigidrot}
\end{equation}
where the rate $\omega\in\R$ is \emph{not} assumed nonzero; that $\omega\ne0$ is a consequence of
nondegeneracy, established as Step~0 of the proof.
Assume also that the repulsive interaction is pairwise \emph{central} along this orbit. (Symmetry of
the interaction need not be assumed: the neighbor set $\calN_i=\{j\ne i:\norm{x_{ij}}\le\mu\}$ is
symmetric by construction, since $j\in\calN_i\iff i\in\calN_j$. Nor is the vanishing of the net
internal torque an assumption; it is derived from centrality in the proof.) Then
\begin{equation}
  \omega^2=k_2,
  \qquad |\omega|=\sqrt{k_2}.
\end{equation}
\end{proposition}

\begin{proof}
Since all pairwise distances are constant under \eqref{eq:rigidrot}, the neighbor graph and the
central force magnitudes are constant along the orbit; hence the repulsive force co-rotates,
$\phi_i(t)=R(\omega t)\phi_i^0$, where $\phi_i^0=\phi_i(\txi^0)$.

\emph{Step 0: $\omega\ne0$.} Suppose $\omega=0$. Then \eqref{eq:rigidrot} reads
$\txi_i(t)\equiv\txi_i^0$, so $\dot{\txi}_i\equiv0$ and $\phi_i(t)\equiv\phi_i^0$, while the observer
equation $\dot{\tvell}_i=-k_2\txi_i^0$ has the constant right-hand side and therefore integrates to
the \emph{affine} function
\begin{equation}
  \tvell_i(t)=\tvell_i(0)-k_2t\,\txi_i^0 .
  \label{eq:omega0_vel}
\end{equation}
Substituting into $\dot{\txi}_i=\phi_i-k_1\txi_i+\tvell_i$ gives, for every $t$ in the interval,
\begin{equation}
  0=\phi_i^0-k_1\txi_i^0+\tvell_i(0)-k_2t\,\txi_i^0 .
\end{equation}
The left-hand side is independent of $t$ and the right-hand side is affine in $t$, so on an interval of
positive length the linear coefficient must vanish: $k_2\txi_i^0=0$ for every $i$. As $k_2>0$ this
forces $\txi_i^0=0$ for all $i$, contradicting $\sum_i\norm{\txi_i^0}^2>0$. Hence $\omega\ne0$, and the
remaining steps may divide by $\omega$.

The observer equation gives
\begin{equation}
  \dot{\tvell}_i=-k_2\txi_i=-k_2R(\omega t)\txi_i^0 .
\end{equation}
Integrating and using $J^{-1}=-J$ gives
\begin{equation}
  \tvell_i(t)=\frac{k_2}{\omega}J\txi_i(t)+c_i,
  \label{eq:vreli_const}
\end{equation}
for some constant vector $c_i$. Substitute \eqref{eq:rigidrot} and \eqref{eq:vreli_const} into
$\dot{\txi}_i=\phi_i-k_1\txi_i+\tvell_i$:
\begin{equation}
  \omega J R(\omega t)\txi_i^0
  =R(\omega t)(\phi_i^0-k_1\txi_i^0+\frac{k_2}{\omega}J\txi_i^0)+c_i .
\end{equation}
The left-hand side and the first term on the right co-rotate with $R(\omega t)$, whereas $c_i$ is
constant. Since $\omega\ne0$ and the identity holds for a continuum of times, the constant term must
vanish: $c_i=0$. Therefore
\begin{equation}
  \tvell_i=\frac{k_2}{\omega}J\txi_i .
  \label{eq:vreli}
\end{equation}
Evaluating the same closed-loop identity at $t=0$ gives
\begin{equation}
  \Bigl(\omega-\frac{k_2}{\omega}\Bigr)J\txi_i^0=\phi_i(\txi^0)-k_1\txi_i^0 .
  \label{eq:torque_eq}
\end{equation}
Take the two-dimensional cross product of $\txi_i^0$ with both sides and sum over $i\in\NN$. Using
$\txi_i^0\times J\txi_i^0=\norm{\txi_i^0}^2$ and $\txi_i^0\times\txi_i^0=0$,
\begin{equation}
  \Bigl(\omega-\frac{k_2}{\omega}\Bigr)\sum_{i\in\NN}\norm{\txi_i^0}^2
  =\sum_{i\in\NN}\txi_i^0\times\phi_i(\txi^0) .
  \label{eq:torque_sum}
\end{equation}
The right-hand side is the net torque of the internal repulsive forces about the target. For each
undirected interacting pair $(i,j)$, the two equal-and-opposite central contributions are parallel to
$\txi_i^0-\txi_j^0$, and their torque contribution is
\begin{equation}
  \txi_i^0\times(\txi_i^0-\txi_j^0)+\txi_j^0\times(\txi_j^0-\txi_i^0)
  =-\txi_i^0\times\txi_j^0+\txi_i^0\times\txi_j^0=0 .
\end{equation}
Thus $\sum_i\txi_i^0\times\phi_i(\txi^0)=0$, and the nondegeneracy assumption gives
\begin{equation}
  \Bigl(\omega-\frac{k_2}{\omega}\Bigr)\sum_{i\in\NN}\norm{\txi_i^0}^2=0
  \quad\Longrightarrow\quad
  \boxed{\;\omega^2=k_2,\qquad |\omega|=\sqrt{k_2}.\;}
  \label{eq:omega_pred}
\end{equation}
\end{proof}

\begin{remark}[Why the trichotomy collapses to two values]\label{rem:trichotomy}
Read as a polynomial identity, \eqref{eq:torque_sum} is $\omega^2=k_2$ only after multiplying through
by $\omega$; the unmultiplied statement $(\omega-k_2/\omega)\sum_i\norm{\txi_i^0}^2=0$ is undefined at
$\omega=0$. It is therefore natural to ask whether admitting $\omega=0$ enlarges the conclusion to the
trichotomy $\omega\in\{0,+\sqrt{k_2},-\sqrt{k_2}\}$. It does not, for the reason isolated in Step~0: a
stationary configuration is incompatible with the observer dynamics unless every vehicle sits exactly
on the target. The mechanism is worth stating separately, because it is the same integral action that
drives the rotation. The observer $\dot{\tvell}_i=-k_2\txi_i$ has no damping term, so a constant
position error $\txi_i^0\ne0$ is integrated without bound and $\tvell_i$ grows linearly in $t$; a
genuinely static formation would need that growing velocity estimate to be balanced by the static
terms $\phi_i^0-k_1\txi_i^0$, which is impossible. Rotation is thus not merely one option among
three---it is the \emph{only} way a persistent nonzero position error can coexist with a
non-damped integral observer.

Two further points make the collapse robust rather than a technicality. First, the excluded case
$\txi_i^0=0$ for all $i$ is not merely excluded by fiat: the collision barrier
$\alpha(s)\to\infty$ as $s\to d^+$ prevents distinct vehicles from coinciding, so vehicles cannot all
occupy the target position. The degenerate root is therefore unreachable for the actual closed loop, not
just ruled out by the nondegeneracy hypothesis. Second, the two surviving roots $\pm\sqrt{k_2}$ are
both genuinely realized, and which one is selected is an initial-condition-dependent
symmetry-breaking question rather than a further consequence of the dynamics; this is the subject of
Section~\ref{sec:direction}. The proposition thus fixes the angular \emph{speed} completely while
leaving the \emph{sign} open, and both signs appear in the sweep of
Section~\ref{S-sec:sim_omega}.
\end{remark}

\begin{remark}[The centroid must sit exactly on the target]\label{rem:centroid}
The ansatz \eqref{eq:rigidrot} in fact forces the average error to vanish, so the assumption
$\bx=0$ made elsewhere is not an extra restriction. Averaging \eqref{eq:torque_eq} over $i\in\NN$ and
using $\sum_i\phi_i=0$ (equal and opposite pairwise forces) gives
$(\omega-k_2/\omega)J\bx^0=-k_1\bx^0$. Equivalently, substituting $\bx(t)=R(\omega t)\bx^0$ into the
average dynamics \eqref{eq:osc} yields
\begin{equation}
  (k_2-\omega^2)\,\bx^0+k_1\omega\,J\bx^0=0 .
\end{equation}
Since $J$ has no real eigenvalue, $\bx^0$ and $J\bx^0$ are linearly independent unless $\bx^0=0$;
as $k_1,\omega\ne0$ the coefficient of $J\bx^0$ cannot be cancelled, so $\bx^0=0$. Hence the
formation centroid of any nondegenerate rigid rotation coincides with the target exactly.
\end{remark}

This is a rigorous \emph{necessary} condition for any nondegenerate rigid rotation. It does not prove
that such a rotation exists, is unique, or attracts nearby trajectories. Those harder questions
require either a nonlinear invariant-manifold/orbital-stability proof or a certified spectral
argument in a co-rotating frame. The important point for the present report is that whenever a rigid
rotation is observed, its angular speed is forced by torque balance and observer integration; no
choice of radius or non-regular shape can change it.

\subsection{Degenerate symmetric subspaces: no rotation}
\label{sec:degenerate}
The theorem is conditional: it fixes the angular speed \emph{if} a nondegenerate rigid rotation
exists. It does not say that every initial condition must leave every symmetry class and start
rotating. Two invariant subspaces show that the caveat has content.

The first is the fully collinear initialization
\begin{equation}
  x_i(0)\in\R\times\{0\},
  \qquad x_0(0)\in\R\times\{0\},
  \qquad v_0\in\R\times\{0\},
  \qquad \bv_i(0)=0 .
  \label{eq:collinear_init}
\end{equation}
Here the $y$-axis reflection symmetry is exact. The attractive term, the observer term, the target
motion, and every central pairwise repulsion have zero $y$ component whenever all agents lie on the $x$
axis, so the subspace \eqref{eq:collinear_init} is invariant, $\txi_{i,2}(t)\equiv0$, and no rotation
can arise. The second is four agents on the coordinate axes with the target fixed at the origin, which
is reflection symmetric across both axes: the forces on the vertical pair stay vertical and those on
the horizontal pair stay horizontal, so again there is no tangential component from which a rotation
could start. Both subspaces instead jam onto the collision barrier, in the manner analyzed in
Section~\ref{sec:jam}.

These are not counterexamples to Proposition~\ref{prop:omega}; they are degenerate invariant subspaces
in which no nonzero rigid rotation is reached. Both are confirmed to machine precision in
Section~\ref{S-sec:sim_degenerate} of the companion.

\subsection{The rotating formation need not be a regular hexagon}
\label{sec:shapebranch}
The radial balance $\phi_i=k_1\txi_i$ that accompanies a rigid rotation does not single out a unique
shape. For the paper's parameters the regular-hexagon branch of that balance solves
$\alpha(r)=k_1r$, i.e.\ $1/(r-5)-1/4=r/2$, giving $r^*\approx5.35$; but the configuration actually
reached from the hexagonal initialization is an interleaved two-triangle with radii alternating
$[3.39,6.31,3.39,6.31,3.39,6.31]$. Several shape branches thus coexist and the initial condition
selects among them.

What is \emph{not} shape-dependent is the rotation rate. All six vehicles share the same angular
velocity $\sqrt{k_2}$ regardless of their differing radii, which is exactly the content of
Proposition~\ref{prop:omega}: the zero-torque identity eliminates every $r_i$ before $\omega$ is
determined. This is the sharpest available illustration that the selection theorem constrains the rate
without constraining the geometry (Section~\ref{S-sec:sim_observations} of the companion).

\subsection{Conjecture: a trichotomy of long-term behaviors}
\label{sec:dist_conjecture}

Proposition~\ref{prop:omega} and Section~\ref{sec:corotating} characterize the angular speed and the
local orbital stability of the rotating branch, and the degenerate cases above show that some
initializations never rotate. Neither settles the harder dynamical question of what long-term
behaviors the system can exhibit. The numerical evidence assembled in the companion document
(Section~\ref{S-sec:sim_observations}) points to three distinct attractor types, which we record
as a conjecture. The jammed alternative is established rigorously for the collinear subspace in
Proposition~\ref{prop:jam} below, and the breathing limit cycle is observed numerically in a
companion parameter sweep (Section~\ref{S-sec:sim_observations}) under different initial conditions
and gain values---most clearly at $k_1=0.5$, $k_2=0.5$, $\mu=9.0$, $d=5.0$ with $N=5$ agents.

\begin{conjecture}[Trichotomy]\label{conj:dist}
Consider the single-integrator closed loop \eqref{eq:veh},\eqref{eq:tgt},\eqref{eq:ctrl7} with the
repulsive magnitude \eqref{eq:alpha} and gains $k_1,k_2>0$. For every initial condition with
$\min_{i\ne j}\norm{x_i(0)-x_j(0)}>d$, the $\omega$-limit set of the shape trajectory
$\txi^s(t)$ is bounded and belongs to exactly one of the following three mutually disjoint
classes of compact invariant sets:
\begin{enumerate}[leftmargin=2.5em]
  \item \emph{Rotating equilibrium orbit:} a nondegenerate rigid rotation about the target at
        angular speed $|\omega|=\sqrt{k_2}$, with all pairwise distances converging to finite
        constants $d_{ij}^*>d$ (the generic case);
  \item \emph{Jammed boundary:} a configuration in which at least one pair attains the collision
        distance $d_{ij}^*=d$, the positions converge, and the observer states $\tvell_i$ diverge
        linearly;
  \item \emph{Breathing limit cycle:} a non-trivial periodic orbit of the shape variables,
        \[\txi^s(t+T)=\txi^s(t)\qquad\forall t\ge T_0,\]
        with all pairwise distances remaining strictly above $d$ and oscillating periodically;
        the angular speed and individual radii vary along the cycle.
\end{enumerate}
In all three cases the average position error converges to zero exponentially
(Section~\ref{sec:avg}), so the formation center locks onto the target.
\end{conjecture}

\paragraph{What is known rigorously.} The average position error converges to zero exponentially
(Section~\ref{sec:avg}); the formation center therefore locks onto the target. The repulsive
magnitude $\alpha(s)\to+\infty$ as $s\to d^+$ forms a logarithmic barrier, so
$\min_{i\ne j}\norm{x_i(t)-x_j(t)}>d$ for all $t$: collision cannot occur for any finite time.
Proposition~\ref{prop:omega} is a necessary condition for the rotating class.

\paragraph{The rotating branch.} For the paper's regular-hexagon initialization the distances converge
to the interleaved two-triangle pattern of Section~\ref{sec:shapebranch}, and the co-rotating
linearization of Section~\ref{sec:corotating} has a single zero mode, so the branch is locally orbitally
stable. Global attraction to this branch is not proved.

\paragraph{The jammed branch.}
\label{sec:jam}
In the collinear subspace this alternative is not conjectural: it can be proved outright, so for that
class the trichotomy collapses to a theorem. The argument below is self-contained and uses only the shape
energy of Section~\ref{sec:corotating}.

\begin{proposition}[Jamming in the collinear subspace]\label{prop:jam}
Consider the single-integrator closed loop restricted to the invariant collinear subspace
\eqref{eq:collinear_init}, with $k_1,k_2>0$ and all initial gaps exceeding $d$. Then the shape
trajectory converges, $\txi^s(t)\to\xi^*$, the limit satisfies $\xi^*\ne0$, and at least one gap
tends to the collision distance $d$. Along the way $\norm{\tvell^s(t)}\to\infty$ linearly.
\end{proposition}

\begin{proof}
In one dimension the repulsive Hessian has only radial directions, with
$\Phi''(s)=\alpha'(s)=-1/(s-d)^2<0$, so $\nabla^2P\preceq0$ and
$B=k_1I-\nabla^2P\succeq k_1I\succ0$. By \eqref{eq:Edot} the shape energy \eqref{eq:energy} satisfies
\begin{equation}
  \dot E\le-k_1\norm{\dot{\txi}^s}^2\le0 ,
  \label{eq:jam_Edot}
\end{equation}
so $E$ is nonincreasing and the shape state is bounded (the barrier bounds it below, coercivity of
the $k_1$ term above). Integrating \eqref{eq:jam_Edot} gives
$k_1\int_0^\infty\norm{\dot{\txi}^s}^2\dd t\le E(0)-\inf E<\infty$, hence
$\dot{\txi}^s\in L^2$; with $\ddot{\txi}^s$ bounded on the bounded shape set this forces
$\dot{\txi}^s\to0$ and therefore $\txi^s\to\xi^*$ for some limit $\xi^*$.

No interior equilibrium exists: stationarity of the observer requires $\dot{\tvell}_i=-k_2\txi_i=0$,
i.e.\ $\txi_i=0$ for every $i$---total collapse onto the target, which the logarithmic barrier
forbids. Hence $\xi^*\ne0$, and consequently $\tvell^s$ cannot be stationary. Explicitly,
$\dot{\tvell}^s=-k_2\txi^s\to-k_2\xi^*\ne0$, so
\begin{equation}
  \tvell^s(t)=-k_2\xi^*\,t+o(t),
  \qquad \norm{\tvell^s(t)}\to\infty
  \label{eq:jam_vdiv}
\end{equation}
at a linear rate. Since $\dot{\txi}^s\to0$ while $\tvell^s$ diverges, the remaining terms of
$\dot{\txi}^s=\phi(\txi^s)-k_1\txi^s+\tvell^s$ must balance it; as $-k_1\txi^s\to-k_1\xi^*$ is
bounded, the repulsion must satisfy $\norm{\phi(\txi^s(t))}\to\infty$. The magnitude $\alpha(s)$ is
finite for every $s>d$ and blows up only as $s\to d^+$, so at least one gap tends to $d$.
\end{proof}

\begin{remark}[Rate of approach to the barrier]
The proof also gives the rate. Let $s(t)$ be a gap tending to $d$ and project the balance above onto
its axis: writing $c^*=k_2\abs{\langle\xi^*,e\rangle}>0$ for the component of the diverging observer
term along that axis $e$, the leading-order cancellation is $\alpha(s(t))\sim c^*t$. With
$\alpha(s)=1/(s-d)-1/(\mu-d)$ this gives
\begin{equation}
  s(t)-d\ \sim\ \frac{1}{c^*\,t}\ \longrightarrow\ 0 ,
  \label{eq:jam_asymp}
\end{equation}
so the gaps close on the barrier at rate $1/t$ while the observer states diverge linearly.
\end{remark}

A stiff (BDF) integration confirms the $1/t$ scaling of \eqref{eq:jam_asymp} to order of magnitude; the
runs are reported in Section~\ref{S-sec:sim_jam} of the companion. Only the scaling is verified, not the
coefficient $c^*$, since \eqref{eq:jam_asymp} is a leading-order balance that discards the $O(1)$
attractive and neighbour contributions.

\paragraph{Why a complete proof is hard.} Proposition~\ref{prop:jam} settles the collinear class
because there the damping matrix is definite. In two dimensions $B=k_1I-\nabla^2P$ is sign-indefinite
away from the rotating equilibrium: the radial directions carry $\alpha'(s)<0$ (damping), while the
tangential stiffness $\alpha(s)/s>0$ contributes anti-damping---the very term that sustains the
rotation. The shape energy is therefore not a global Lyapunov function, and no monotone quantity
separating the rotating and jammed attractors is known. The jammed branch, moreover, lies on the
singular boundary $s=d$ where the vector field is undefined, so standard convergence theorems do not
apply.

\paragraph{Numerical caveat.} The barrier makes the problem stiff: near $s=d$ the local timescale is
$s-d\to0$ while the stiffness $|\alpha'(s)|=1/(s-d)^2\to\infty$. Fixed-step explicit schemes therefore
blow up near the barrier and can misleadingly suggest divergence or large-amplitude oscillation; a
stiff solver is required to observe the bounded, converging dynamics. This is why the jammed-branch
runs in the companion use BDF rather than the RK4 scheme used elsewhere.

